% =====================================================================
\section{Threats to validity}
\label{sec:threats}
% =====================================================================

\paragraph{Construct validity}
The plant is our own model, and the planner plans with that same model:
it is an oracle of the system it controls, which flatters \jcac{} against
comparators that see only utilisation or request rates. On a plant
re-fitted to the deployed gateway's measurements, a planner that kept the
published model lost to the layered stack by $0.1052$ in $J$. The plant's
replicas are gateway pods, not GPU model servers: batching, KV-cache
memory, request backlog and tier-switch lag lie outside the study, and
pod congestion multiplies the external model's latency instead of adding a
queueing delay. A replica-step costs about one small-tier cache miss while
a tenant sends tens of misses per step, so the bill is mostly model spend,
set by a synthetic price ladder of 1:10:100 per miss in ten of the eleven
campaigns; our own GPU bench measured 1:1.475:1.518. Under the measured
ladder coordination was worth nothing (QUALITY), and the complete-policy
verdicts move: in an exploratory re-run of JOINTNESS under that ladder,
\jcac{} beat every comparator, the cheapest-tier stack by $0.0276$
$[0.0045,0.0531]$, but only because the comparators stayed pinned to the
small tier and because small-tier agent requests are slow by assumption.
That assumption (6\,s on the small tier against 2.5 and 3.5\,s on the
larger ones) is unmeasured and contradicted by our own bench; with a
monotone ladder the directional coordination gain disappeared and
\jcac{}'s wins over capped KEDA and the static share became ties
(exploratory). At the level field of LEVEL the monotone row changed no
verdict (LV-R1), so the final RQ1 answer does not rest on it; the
synthetic price ladder was not varied there and remains an assumption under
every final verdict. RQ2's conclusion that coordination is small holds under
both rows and both ladders. The cache is a synthetic hit
curve, bounded by the CS campaign but with no model of similarity,
eviction, drift or wrong answers; its memory is almost free (0.2\% of the
bill), so more cache is always better up to the caps. $J$ uses the planner's own
weights; it ignores budget overruns (3.1--10.4\% of the comparators'
tenant-steps), prices a shed request only through a capped violation term,
leaves network, storage and control-plane cost unpriced, and includes
answer quality only in QUALITY. In an exploratory re-scoring, the registered
violation weight of 2 was exactly the crossover against re-tuned HPA (at 4
HPA was better, at 8 KEDA too) and the re-scoring did not re-optimise the
arms at each weight. The cost per request served changes sign with its
aggregation (Section~\ref{sec:shed}), so we make no per-request cost claim.
The budgets were set below the small tier's peak spend for several
workloads, so part of the shedding is a property of the budgets we chose.

\paragraph{Internal validity}
The comparators are hand-written rules with one or two parameters, tuned
on a 12-run slice; per-size re-tuning moved their cost and violation but
barely their $J$, because $J$ is nearly flat in their targets. The
registered layered stacks were not held to the cluster caps, and the
cheapest-tier stack used up to twice the small cluster's replicas; its
capped re-run is exploratory and was not re-tuned under the caps. In the
first eleven campaigns only \jcac{} could shed or keep a budget, and the
budget-matched campaign could not remove that asymmetry
(Section~\ref{sec:budget}). LEVEL removed both leaks (every arm capped and
re-tuned; a budget-gated stack), but its own choices carry risk: the
treatment variant was selected on the comparators' tuning slice under
instant actuation, which is how a variant fragile under lag was chosen; the
gate sheds whole steps, as \jcac{} does, rather than refusing only the
requests beyond the budget; and the comparators were not re-tuned for lag,
held-out demand or the trace. \jcac{} plans one held move
per interval on a three-point trend forecast; a damped Holt forecaster cut
one-step error on the BurstGPT trace by 6.9--26.9\%, but joint and
decoupled search share the forecast, so RQ2 is unaffected. A multi-step
planner might find sequenced coordination that a held move cannot; the
wide-lattice \jcac{} beat neither the default-behaviour nor the capped HPA
(RV-W1, RV-W2). The solver is joint within a tenant and coordinate descent
across tenants; on the earlier, non-convergent solver, six sweep orders moved
mean Jain fairness by at most 0.0076 and cost by 0.91\%. Actuation is
instantaneous outside ACTUATION, and even there tier switches never lag.
Besides the comparator defects of Section~\ref{sec:chain}, our audits found
two defects that hurt \jcac{} (the non-convergent solver, replaced in every
campaign here, and the comparators' missing budget, which BM measures) and
several in the live harness (Section~\ref{sec:live}).

\paragraph{External validity}
The main workloads are synthetic and periodic, every design point puts
eight tenants of a single workload class together (so tenants interact
only through the cluster caps), HELDOUT has one AI class, the coupled plant was chosen from a probe grid, and the live-fitted
plant is a parameter scenario, not a live experiment: no measured run
reached congestion and AI requests never loaded a replica, so the coupling
that coordination needs was absent from the deployed system's
measurements. The two non-periodic workloads and the BurstGPT windows,
where \jcac{} beat no tuned comparator, are the closest we come to
deployment conditions. Nothing in Sections~\ref{sec:chain}--\ref{sec:width}
ran on a cluster, and the live case study of Section~\ref{sec:live} is one
host, two replications and an in-sample calibration. Our conclusions are
conditional on the simulated workloads and coupling.

\paragraph{Statistical conclusion validity}
The intervals resample design points, so they generalise only over the
60 cells we chose; with the five workload classes as the unit, the
confirmatory tests we re-checked had $p$-values between 0.25 and 0.997. Replications are
nearly redundant (intraclass correlation 0.97--0.99), so the effective
sample is the design point, not the run. The $\pm0.025$ margin is the
design's minimum detectable difference, not an operational threshold; an
operator who cares about 1\% of $J$ should read Table~\ref{tab:coord} as a
small directional effect. Several registered passes do not survive a
programme-wide correction (\ref{app:tests}). The programme is sequential
and adaptive: most campaigns were designed after an exploratory reading of
the same 60 design points, so their confirmatory status protects against
analytic flexibility but not against the choice of what to test next; only
HELDOUT tested design points no earlier analysis had seen, and LEVEL re-ran
the known design points on fresh seeds.

\paragraph{Registration evidence}
A registration is only as strong as its timestamp. The protocols of the
seven core campaigns and of LEVEL were pushed to the public repository
before their first runs (LEVEL's 15\,s before its runner started, with the
push recorded by a CI run); those of RV, TA, CS, BM and the SA re-run were
committed locally before their runs and published afterwards, and six
protocol--campaign pairs began within two minutes of their protocol's
commit. LEVEL's scorer needed one post-run loader fix, disclosed in a dated
note, which changed no number. Software
Heritage holds pre-run snapshots for two protocols only, and the Open
Science Framework registration is a single post-run mirror of the first 48
protocols. The public history was later rewritten twice to edit commit
messages and a few non-protocol files; no protocol, record, data file or
figure changed, commit timestamps were preserved, and a map in the artifact
resolves the old commit identifiers that server-side witnesses cite.
